\documentclass{beamer}
\usetheme{Madrid}
\usecolortheme{default}

% UdeA Color Setup
\definecolor{UdeAGreen}{RGB}{0, 115, 54}
\setbeamercolor{structure}{fg=UdeAGreen}

\title[Orbit Determination]{Fundamentals of Orbit Determination}
\subtitle{Sequential Estimation with Kalman Filters}
\author[Prof. Juan]{Prof. Juan F. Puerta I.}
\institute[UdeA]{Universidad de Antioquia}
\date{\today}

\begin{document}

\begin{frame}
    \titlepage
\end{frame}

\section{Introduction}

\begin{frame}{What is Orbit Determination (OD)?}
    \textbf{Definition:} The process of computing a spacecraft's precise state (position and velocity) based on observational data that is inherently noisy and incomplete.
    
    \vspace{0.5cm}
    \begin{itemize}
        \item \textbf{Input:} Noisy sensor measurements (Range, Doppler, Angles).
        \item \textbf{Goal:} Reconstruct the 3D trajectory $[x, y, z, \dot{x}, \dot{y}, \dot{z}]$.
        \item \textbf{Challenge:} Sensors are imprecise and physical models are non-linear.
    \end{itemize}
\end{frame}

\section{Why Kalman Filters?}

\begin{frame}{The Need for Kalman Filters}
    Classical methods (Batch Least Squares) are excellent for ground-based post-processing, but they are limited for modern autonomous navigation.
    
    \begin{block}{Why KFs are Essential}
        \begin{enumerate}
            \item \textbf{Real-Time Processing:} Recursive structure requires only the current state and measurement—no massive historical data storage needed.
            \item \textbf{Sequential Updates:} Updates the orbit the moment a new ping arrives, crucial for autonomous station-keeping.
            \item \textbf{Stochastic Balancing:} Mathematically balances "trust" between the physics model ($Q$) and sensor accuracy ($R$).
        \end{enumerate}
    \end{block}
\end{frame}

\section{The Filter Family}

\begin{frame}{The Kalman Filter Family}
    \begin{itemize}
        \item \textbf{KF (Linear):} Best Linear Unbiased Estimator (BLUE). Optimal for linear systems. \textit{Rarely used in orbit determination due to non-linearity.}
        \item \textbf{EKF (Extended):} The industry workhorse. Uses \textbf{Jacobians} (1st order Taylor expansion) to linearize non-linear orbit equations at each step.
        \item \textbf{UKF (Unscented):} Uses \textbf{Sigma Points} to capture non-linearities to higher order. Superior for high-tumbling or highly perturbed trajectories.
    \end{itemize}
\end{frame}

\section{The EKF Mechanism}

\begin{frame}{The Predict-Correct Cycle}
    \begin{columns}
        \begin{column}{0.5\textwidth}
            \textbf{1. Predict Step}
            \begin{itemize}
                \item Propagate physics model: $\hat{\mathbf{x}}_k^- = \int f(\hat{\mathbf{x}}_{k-1}^+) dt$
                \item Propagate uncertainty: $P_k^- = \Phi P_{k-1}^+ \Phi^T + Q$
            \end{itemize}
        \end{column}
        \begin{column}{0.5\textwidth}
            \textbf{2. Correct Step}
            \begin{itemize}
                \item Calculate Residual: $y_k = z_k - h(\hat{\mathbf{x}}_k^-)$
                \item Update State: $\hat{\mathbf{x}}_k^+ = \hat{\mathbf{x}}_k^- + K_k y_k$
                \item Minimize Covariance: $P_k^+$
            \end{itemize}
        \end{column}
    \end{columns}
\end{frame}

\section{Conclusion}

\begin{frame}{Key Takeaway for Students}
    \begin{quote}
        "Propagating an orbit is looking forward: 'If I know exactly where the satellite is now, where will it be tomorrow?' Orbit Determination is looking backward: 'Based on these noisy sensor pings from the last three days, where must the satellite actually be right now?'"
    \end{quote}
    
    \vspace{0.5cm}
    \centering \textbf{The Kalman Filter is the bridge between chaotic real-world sensor data and precise orbital physics.}
\end{frame}

\end{document}